% Appendix D -- two design figures moved out of section 4 (2026-09-26) for the
% page limit. Both illustrate structure the body states in prose; the three that
% stayed (the per-layer execution split, the compaction, and the k-slot
% pipeline) show things prose cannot.

\section{Additional design figures}
\label{app:figures}

The data path is Fig.~\ref{fig:arch} in \S\ref{sec:design} --- which component talks to which,
and where the host mirror sits relative to the device pool.
Figure~\ref{fig:lifecycle} gives a request's states and the transitions between
them. Both illustrate structure \S\ref{sec:design} states in prose.

% (fig:arch moved to the body; see sections/fig-arch.tex)


\begin{figure}[t]
\centering
\resizebox{\columnwidth}{!}{%
\begin{tikzpicture}[
  x=1cm, y=1cm, font=\scriptsize,
  st/.style={draw, rounded corners=2pt, align=center, minimum height=7.5mm,
             minimum width=1.95cm, inner sep=1.2mm, fill=black!3},
  hs/.style={st, fill=blue!10, draw=blue!45},
  arr/.style={-{Latex[length=1.5mm]}},
  ed/.style={font=\scriptsize\itshape, text=black!60, align=center},
]
\node[st] (wait) at (0,0) {waiting};
\node[hs] (adm)  at (2.9,1.35) {admitted host-tier\\{\scriptsize 0 device blocks}};
\node[hs] (dec)  at (6.0,1.35) {decoding\\host-tier};
\node[st] (res)  at (6.0,-1.35) {decoding\\device-tier};
\node[st] (fin)  at (9.6,0) {finished};

\draw[arr] (wait) -- node[ed, above, sloped] {tier $=$ host} (adm);
\draw[arr] (wait) -- node[ed, below, sloped] {tier $=$ device} (res);
\draw[arr] (adm)  -- node[ed, above, align=center] {\textbf{worker ack}\\KV present} (dec);
\draw[arr] (dec)  -- node[ed, above, sloped] {done} (fin);
\draw[arr] (res)  -- node[ed, below, sloped] {done} (fin);

% migration: two vertical arrows, separated so their labels never meet
\draw[arr] ([xshift=-7mm]res.north) -- ([xshift=-7mm]dec.south);
\node[ed, anchor=east, align=right] at (5.15,0)
     {\textbf{demote}\\on seed ack:\\release blocks};
\draw[arr, dashed] ([xshift=7mm]dec.south) -- ([xshift=7mm]res.north);
\node[ed, anchor=west, align=left] at (6.85,0)
     {\textbf{promote}\\atomic refill};
\end{tikzpicture}}
\caption{\textbf{Request lifecycle.} A request admitted directly to the host
tier never allocates a device block. A resident request that is demoted releases
its blocks only after the worker acknowledges that its KV reached the host ---
releasing on the assumption that it did would destroy the request's history.
Promotion refills the block list atomically; a partial refill is treated as a
failed promotion and the request stays on the host tier.}
\label{fig:lifecycle}
\end{figure}

\begin{figure}[t]
\centering
\begin{tikzpicture}[
  x=1cm, y=1cm, font=\scriptsize,
  blk/.style={draw, minimum width=6mm, minimum height=5mm, inner sep=0pt},
  used/.style={blk, fill=blue!12, draw=blue!50},
  free/.style={blk, fill=black!3, draw=black!20},
  dst/.style={blk, fill=blue!25, draw=blue!60},
  lbl/.style={font=\scriptsize\itshape, text=black!60, align=left},
  arr/.style={-{Latex[length=1.4mm]}, blue!55},
]
\node[font=\bfseries\scriptsize, anchor=west] at (0,1.75) {host mirror, addressed by host block id};
\foreach \i in {0,1,4,5,7} \node[free] at (\i*0.63,1.15) {};
\foreach \i/\t in {2/7,3/8,6/23,8/40,9/41,10/42}
  \node[used] (h\i) at (\i*0.63,1.15) {\tiny \t};

\node[lbl, anchor=west] at (0,0.62)
  {A:\,$\{7,8\}$\quad B:\,$\{23\}$\quad C:\,$\{40,41,42\}$};

\node[font=\bfseries\scriptsize, anchor=west] at (0,-0.55) {rolling slot $\ell \bmod k$ \ ($R{=}6$, 3 runs)};
\foreach \i in {0,...,5} \node[dst] (d\i) at (\i*0.63,-1.2) {\tiny \i};

\draw[arr] (h2.south) -- (d0.north);
\draw[arr] (h3.south) -- (d1.north);
\draw[arr] (h6.south) -- (d2.north);
\draw[arr] (h8.south) -- (d3.north);
\draw[arr] (h9.south) -- (d4.north);
\draw[arr] (h10.south) -- (d5.north);

\node[anchor=north west, font=\scriptsize] at (4.5,-0.45)
  {\begin{tabular}{@{}l@{\ }l@{}}
     \multicolumn{2}{@{}l}{\bfseries compacted block table}\\
     A & $[0,\,1]$\\
     B & $[2]$\\
     C & $[3,\,4,\,5]$\\
   \end{tabular}};
\end{tikzpicture}
\caption{\textbf{Compaction, once per step.} Three host-tier requests hold
scattered host blocks. Their union is sorted and mapped onto $[0,R)$, giving
both the rolling region's size for this step and a block table expressed in
rolling indices. The sorted set decomposes into three maximal \emph{runs}, so
the per-layer transfer is three contiguous slice copies rather than six
scattered ones; with the mirror laid out block-major, each run is one contiguous
byte range on both sides, so the DMA engine performs the compaction and no
host-side gather is required.}
\label{fig:compact}
\end{figure}

% (figure moved to the body; see sections/fig-*.tex)


% Moved here from section 4 (2026-09-27): a six-symbol table, three of whose
% symbols section 3 already defines, is not worth a float in a 10pp body.
\begin{table}[t]
\centering
\small
\begin{tabular}{ll}
\toprule
symbol & meaning \\
\midrule
$L$ & layers in the model \\
$b$ & tokens per KV block (vLLM block size) \\
$k$ & rolling slots, $k \ll L$ (we use 4) \\
$R$ & blocks the \emph{active} host-tier cohort occupies this step \\
$S$ & context length \\
$\btok$ & KV bytes per token, $= 2 L\,n_{kv}\,d_h \cdot \mathrm{sizeof(dtype)}$ \\
\bottomrule
\end{tabular}
\caption{Notation local to the design section. Symbols shared with \S\ref{sec:model} ($S$, $\btok$, $L$) carry the same meaning there; $b$, $k$ and $R$ appear only here.}
\label{tab:notation}
\end{table}

% fig:exec moved here from section 4 (2026-09-27). The prose states what it
% shows; the drawing is the confirmation, not the explanation.
% (figure moved to the body; see sections/fig-*.tex)

